\begin{abstract}
Accurate mortality risk stratification from electronic health records (EHR)
requires models that reason over temporal disease trajectories, numeric
biomarkers, and demographic context.
It remains unclear whether large language models (LLMs) can encode the
\emph{structural} properties of clinical data, including numeric precision,
temporal ordering, and causal relationships, required for calibrated
survival prediction.

We present a controlled three-axis study on a UK Biobank respiratory
disease subgroup ($n\,=\,124{,}158$; 18{,}623 test), a multimorbid
population where the temporal complexity of disease trajectories makes
structural encoding choices particularly consequential.
The study systematically varies:
(A) \emph{input representation} across five temporal-encoding settings
(cross-sectional, sequential ordinal, isolated prose, decoupled time and
event, and unified summary prose);
(B) \emph{encoder architecture} (Qwen3-8B vs.\ Bio\_ClinicalBERT) for the
three embedding-based settings; and
(C) \emph{feature fusion method} (concatenation vs.\ element-wise
summation) for settings that combine embeddings with raw features.
All embedding variants use a uniform Cox proportional hazards head,
enabling direct comparison of structural encoding choices.

Three key findings emerge from our 12-cell experiment matrix.
First, \emph{temporal ordering is the critical structural property}:
decoupled age--event encoding (Setting 4) achieves the best Harrell
C-index (0.624, $+$0.005 over the CoxPH baseline and $+$0.008 over prose
serialisation of the same events).
Second, \emph{LLMs can encode numeric biomarker structure from prose}: a
Qwen3-8B embedding of full clinical prose with no raw numeric features
matches the 39-feature CoxPH baseline (C\,=\,0.619 vs.\ 0.619) and
achieves superior TD-AUC.
Third, \emph{autoregressive generation has a structural calibration
gap}: Delphi shows poor near-term AUC (0.52 at 9.5\,yr) and substantially
worse calibration (IBS\,=\,0.185 vs.\ 0.141) despite competitive long-horizon
discrimination.
These results establish a reproducible benchmark and suggest that
\emph{how} temporal structure is encoded matters more than which specific
LLM encodes it, a claim supported by Bio\_ClinicalBERT producing near-identical
C-indices to Qwen3-8B across Settings 3 and 4.
\keywords{survival analysis \and LLM embeddings \and EHR \and
          UK Biobank \and structured multimodal intelligence}
\end{abstract}
